\section{A sparse macrograph for End-of-Line}\label{sec:macrograph}

We first replace the circuit transitions by a graph in which one edge
changes at most one tape bit and its occupation bit.  The graph is
specified by short labels and local transition rules; its vertices are
not enumerated.  The use of a clock to expose sparse circuit updates
has a close precedent in Rubinstein's LocalEndOfALine
construction~\cite[Sections~5--5.1]{Rubinstein2016Settling}, where a
counter schedules activation and deactivation of individual wires.
We give the transcript, reciprocal proposals, and bounded-incidence
validity description needed here explicitly.

Throughout the outer construction, let $n_0$ be the
number of input bits of the End-of-Line instance and let $N$ be its
complete binary encoding length.  We use the promises
\begin{equation}\label{macro:promises}
 P(0^{n_0})=0^{n_0},\qquad S(0^{n_0})\ne0^{n_0},\qquad
 P(S(0^{n_0}))=0^{n_0}.
\end{equation}
Thus a nonzero string $y$ is an answer if
\begin{equation}\label{macro:answer}
 P(S(y))\ne y\quad\hbox{or}\quad S(P(y))\ne y.
\end{equation}

\subsection{Sparse transcripts and their two orders}\label{macro:transcripts}

Normalize both circuits to gates of fan-in at most two.  In particular,
replace any larger AND or OR by a binary tree.  Give every input and
every gate a distinct tape position.  For each listed output, introduce
a separate terminal copy gate, even if several listed outputs originally
name the same wire.  Replace the outgoing uses of a wire by a binary
tree of copy gates, including its uses as feeds to terminal copies.
A wire with $d$ uses needs $O(d)$ such gates.  The total number of
positions is therefore bounded by a constant times the number of gates,
connections, and input and output list entries in the original instance,
and hence is $O(N)$.

These changes preserve $S$ and $P$.  They also make every tape position
an operand of only a bounded number of internal gate rules.  Choose a
topological order of the nonterminal gates of each circuit, followed by
all its terminal output copies.  Such an order exists because those
copies have no internal consumers: any branching needed for other uses
was placed before them.  Copy gates here are auxiliary transcript rules
and need not belong to the source circuit's original basis.

Pad the tape with dummy positions to a power of two $h=2^\ell$.
An arbitrary fixed lower cutoff on $h$ may be imposed, and still
$h=O(N)$.  Write $[h]=\{0,\ldots,h-1\}$.  On these same positions define
the two permutations
\begin{align*}
 \pi_{\mathrm f}:\;&S\text{ inputs},\ S\text{ gates},\
 P\text{ inputs},\ P\text{ gates},\ \text{dummies},\\
 \pi_{\mathrm r}:\;&P\text{ inputs},\ P\text{ gates},\
 S\text{ inputs},\ S\text{ gates},\ \text{dummies},
\end{align*}
using the chosen gate order inside each block.  For a content array
$D\in\{0,1\}^{[h]}$, let $\operatorname{pref}_{\pi,k}(D)$ agree with
$D$ on $\pi(0),\ldots,\pi(k-1)$ and be zero elsewhere.  Its occupation
array is the indicator of these $k$ positions.  Empty and full tapes
have occupation arrays $\mathbf 0$ and $\mathbf 1$, respectively.
The content and occupation arrays are separate: an occupied position
may contain zero.

For a full content array $D$, the notation
$\operatorname{in}_S(D)$, $\operatorname{out}_S(D)$,
$\operatorname{in}_P(D)$, and $\operatorname{out}_P(D)$ denotes the
vectors at the designated input or terminal output positions.  A
\emph{valid full transcript} $Q_x$ computes $S$ on input $x$, copies
its output to the inputs of $P$, computes $P$, and satisfies
\begin{equation}\label{macro:full-transcript}
 \operatorname{in}_S(Q_x)=x,\qquad
 \operatorname{in}_P(Q_x)=\operatorname{out}_S(Q_x),\qquad
 \operatorname{out}_P(Q_x)=x.
\end{equation}
Its dummy positions are zero.  Consequently $Q_x$ exists precisely
when $P(S(x))=x$, and is unique when it exists.

For every $y\in\{0,1\}^{n_0}$, define two complete computation templates,
without imposing any terminal comparisons.  The template $U_y$ first
computes $S(y)$ and then $P(S(y))$; the template $V_y$ first computes
$P(y)$ and then $S(P(y))$.  Both occupy the same tape coordinates and
have zero dummies.  In $V_y$, the $P$ inputs are $y$ and the $S$ inputs
are $P(y)$.  Thus $U_y$ is a valid full transcript exactly when
$P(S(y))=y$, in which case it is $Q_y$; and $V_y$ is a valid full
transcript exactly when $S(P(y))=y$, in which case it is $Q_{P(y)}$.

\begin{definition}[Computed prefixes]\label{macro:prefixes}
A \emph{forward prefix of length $k$ with parameter $y$} has the first
$k$ positions in $\pi_{\mathrm f}$ occupied.  At occupied positions it
obeys these rules: $S$ inputs copy the corresponding bits of $y$;
$P$ inputs copy the corresponding designated $S$ outputs; every gate
obeys its Boolean or copy rule; and dummies are zero.  In addition,
each occupied designated $P$ output equals the corresponding bit of
$y$.  Unoccupied contents are zero.

A \emph{reconstruction prefix of length $k$ from $y$} uses
$\pi_{\mathrm r}$.  Its occupied $P$ inputs copy $y$, its occupied
$S$ inputs copy the designated $P$ outputs, and its occupied gates
and dummies obey the same rules as above.  Each occupied designated
$S$ output must additionally equal the corresponding bit of $y$.
Unoccupied contents are again zero.
\end{definition}

Every operand in a prefix rule is either an earlier occupied position
or a specified bit of $y$.  In particular, a designated comparison is
an extra test on the newly computed terminal copy, not a substitute
for its computation rule.

\begin{lemma}[Unique prefixes and their first obstruction]
\label{macro:prefix-uniqueness}
For a fixed $y$, the valid forward prefixes exist uniquely for an
initial interval of lengths $0\le k\le k_{\mathrm f}(y)$ and equal
$\operatorname{pref}_{\pi_{\mathrm f},k}(U_y)$.  The valid
reconstruction prefixes likewise exist uniquely for
$0\le k\le k_{\mathrm r}(y)$ and equal
$\operatorname{pref}_{\pi_{\mathrm r},k}(V_y)$.  Moreover,
\begin{align}
 k_{\mathrm f}(y)=h&\quad\Longleftrightarrow\quad P(S(y))=y,
 \label{macro:forward-obstruction}\\
 k_{\mathrm r}(y)=h&\quad\Longleftrightarrow\quad S(P(y))=y.
 \label{macro:reverse-obstruction}
\end{align}
If a maximal length is less than $h$, the next position is a designated
output at which the additional comparison fails.
\end{lemma}

\begin{proof}
Induct along the relevant order.  Input copies, gates, and dummies
uniquely determine the next content bit from data already determined.
These are exactly the bits of $U_y$ or $V_y$.  The only possible
obstruction is an additional terminal comparison.  Shortening a valid
prefix removes requirements without changing retained bits, so the
valid lengths form an initial interval.  At length $h$, all terminal
comparisons have been imposed, giving the two equivalences.
\end{proof}

\subsection{The cyclic schedule and reciprocal proposals}
\label{macro:schedule}

A logical tape vector is
\[
 b=(A,E^A,B,E^B)\in\{0,1\}^{4h},
\]
where $A,B$ are contents and $E^A,E^B$ are occupations.  The schedule
uses four phases.  To see their effect when all computations succeed,
suppose $Q_x$ and $Q_y$ exist and $y=S(x)$.  Starting with $A=Q_x$
and $B$ empty, it fills $B$ with $Q_y$, erases $A$ by reversing
reconstruction from $y$, copies $B$ into $A$, and erases $B$.
The resulting tapes are $A=Q_y$ and $B$ empty.  The identity
$P(y)=x$, which follows from validity of $Q_x$, is what makes the
reconstruction recover the old tape.

For intermediate prefixes, including those at a failed terminal
comparison, validity is defined as follows.  The stage $s$ lies in
$\mathbb Z/(4h)\mathbb Z$.  Write $\mathcal V_s$ for the following
set of valid logical vectors; the endpoint descriptions in adjacent
cases are identified.
\begin{enumerate}[label=\textup{(\roman*)}]
 \item $0\le s\le h$: $A=Q_x$ is full for some $x$, and $B$ is a
 forward prefix of length $s$ with parameter
 $y=\operatorname{out}_S(A)=S(x)$.
 \item $h\le s\le2h$: $B=Q_y$ is full for some $y$, and $A$ is a
 reconstruction prefix of length $2h-s$ from
 $y=\operatorname{in}_S(B)$.
 \item $2h\le s\le3h$: $B$ is a valid full transcript and
 $A=\operatorname{pref}_{\pi_{\mathrm f},s-2h}(B)$.
 This copied prefix has no additional computation requirements.
 \item $3h\le s\le4h$: $A$ is a valid full transcript and
 $B=\operatorname{pref}_{\pi_{\mathrm f},4h-s}(A)$.
\end{enumerate}
Occupations in every case are the stated full or prefix indicators.
Stage $4h$ is stage $0$, at which $A$ is full and $B$ is empty.

\begin{lemma}[Boundary agreement]\label{macro:boundaries}
The two validity descriptions agree at each of $h,2h,3h$, and $4h=0$.
\end{lemma}

\begin{proof}
At stage $h$, the first description gives
\[
 A=Q_x,\qquad B=Q_y,\qquad y=S(x),\qquad P(y)=x.
\]
The reconstruction template from $y$ computes $P(y)=x$ and then
$S(x)=y$, so it agrees with $A$ at every position and satisfies all
its comparisons.  This gives the second description.  Conversely,
the second description at $h$ gives $B=Q_y$ and a full reconstruction
$A=V_y$ with $S(P(y))=y$.  Put $x=P(y)$.  Then
$P(S(x))=P(y)=x$, so $A=Q_x$ and $\operatorname{out}_S(A)=y$.
Since $B=Q_y$, it is the full forward prefix required in the first
description.  These identities include all copy gates and dummies.
At $2h$, both descriptions have $A$ empty and $B$ a valid full
transcript.  At $3h$, both have $A=B$ a valid full transcript.
At $4h=0$, both have $A$ full and $B$ empty.
\end{proof}

For each \emph{start frame} $c\in\{0,\ldots,4h-1\}$, define a
forward proposal $F_c$ from stage $c$ to stage $c+1$, and a reverse
proposal $R_c$ from stage $c+1$ to stage $c$, with stages reduced
modulo $4h$.  Thus at current stage $s$, the successor uses $F_s$
and the predecessor uses $R_{s-1}$, where the latter subscript is
reduced modulo $4h$.  The following table specifies their content changes.
Every fill or copy sets the corresponding occupation to one, and
every erasure sets it to zero.  All other bits remain unchanged.
\begin{center}
\begin{tabular}{@{}llll@{}}
\toprule
Start frame & Position $j$ & $F_c$ & $R_c$ \\
\midrule
$0\le c<h$ & $\pi_{\mathrm f}(c)$ & fill $B_j$ forward & erase $B_j$ \\
$h\le c<2h$ & $\pi_{\mathrm r}(2h-c-1)$ & erase $A_j$ & fill $A_j$ by reconstruction \\
$2h\le c<3h$ & $\pi_{\mathrm f}(c-2h)$ & copy $B_j$ to $A_j$ & erase $A_j$ \\
$3h\le c<4h$ & $\pi_{\mathrm f}(4h-c-1)$ & erase $B_j$ & copy $A_j$ to $B_j$ \\
\bottomrule
\end{tabular}
\end{center}
A computed fill always writes its Boolean rule value, even if the
additional terminal comparison fails.  Thus the proposal itself is
defined separately from the test of whether its result is valid.

The \emph{macrograph} has vertex set
\[
 \mathcal V=\{(s,b):s\in\mathbb Z/(4h)\mathbb Z, b\in\mathcal V_s\}.
\]
It has an edge $(c,b)\longrightarrow(c+1,b')$ precisely when
$b'=F_c(b)$ and both endpoints are valid.

\begin{lemma}[Reciprocity and local existence tests]
\label{macro:local-edges}
Whenever $b\in\mathcal V_c$ and $b'\in\mathcal V_{c+1}$,
\begin{equation}\label{macro:reciprocity}
 b'=F_c(b)\quad\Longleftrightarrow\quad b=R_c(b').
\end{equation}
For a valid input to either proposal, validity of its output is
determined by a bounded number of bits at public, frame-dependent
addresses.  The forward proposal can fail only when $0\le c<h$ and
the new forward terminal comparison fails.  The reverse proposal can
fail only when $h\le c<2h$ and the new reconstruction terminal
comparison fails.  The graph has indegree and outdegree at most one,
and no self-edges.
\end{lemma}

\begin{proof}
In the first range, erasing the newly filled position recovers its
previous blank content and occupation.  Conversely, a valid longer
forward prefix has the uniquely prescribed value at its last
position, so filling its shortened prefix recovers it.  In the second
range, erasure shortens a reconstruction prefix.  Its erased value
is uniquely determined by earlier retained positions and the
unchanged full tape $B$, so reverse reconstruction recovers it.
The converse follows from the same uniqueness.  In the last two
ranges, the unchanged full tape determines the copied value,
and erasure recovers the blank position.  Boundary agreement applies
to each use of these statements at a boundary.  In particular,
at $c=4h-1$, the reverse proposal from stage zero copies
$A_{\pi_{\mathrm f}(0)}$ into the first forward position of the empty
$B$ tape, exactly reversing the wraparound erasure.

Shortening a valid computed prefix is always valid, as are extending
or shortening a copied prefix.  When extending a computed prefix,
all earlier conditions remain true.  Its next value is obtained
from at most two earlier operands or a specified full-tape input
bit, so the only new condition that can fail is the stated terminal
comparison.  That comparison needs one additional specified bit.
Thus the relevant operands, old bits, and comparison bits suffice
to decide both proposals and the corresponding existence flags.
Equation~\eqref{macro:reciprocity} gives at most one predecessor,
and the definition gives at most one successor.  Every edge changes
the stage by one modulo $4h$, and $4h\ge4$, so no edge is a self-edge.
\end{proof}

\subsection{Endpoints, sparse edits, and uniform size}
\label{macro:interface}

The following statement is the interface used by the evolving
representation.

\begin{lemma}[Macrograph interface]\label{lem:macrograph}
From every promised End-of-Line instance of explicit length $N$, one
can deterministically construct, in time polynomial in $N$, the
above tape specification with $h=2^\ell=O(N)$, the stage conditions
$\mathcal V_c$, and the following data and algorithms.
\begin{enumerate}[label=\textup{(\roman*)}]
 \item The macrograph has indegree and outdegree at most one, no
 self-edges, and no isolated vertices.  It has a distinguished source
 $v_*$ whose label and actual successor are computable in polynomial
 time.  Every endpoint other than $v_*$ supplies a nonzero
 End-of-Line answer by reading a designated vector on a full tape.
 \item Forward and reverse proposals at the same start frame are
 reciprocal on valid pairs.  Their edits and existence flags use
 only a fixed number of logical bits at public addresses, with a
 fixed finite set of Boolean decision rules.
 \item There is a fixed finite alphabet
 $\Sigma=\{-1,0,1\}^4$ and public edit vectors $h_\sigma(c)$ such
 that each chosen proposal has the form
 $b\mapsto b+h_\sigma(c)$.  One public address is allotted to each
 logical array.  On actual edges only the active tape's content
 position and its occupation position can change.  The same fixed
 edit formulas are defined over any field, whether or not their
 arguments represent valid Boolean tapes.
 \item For each fixed phase, the local tape requirements have
 bounded arity and bounded operand multiplicity, uniformly in $N$.
 The full tape specification and all frame-dependent rule and
 address tables have $O(h)$ records of $O(\log(N+2))$ bits each.
 A vertex label has $4h+\lceil\log_2(4h)\rceil=O(N)$ bits;
 validity and endpoint decoding take polynomial time.  These
 bounds concern descriptions and individual labels and do not
 require enumeration of the macrograph.
\end{enumerate}
\end{lemma}

\begin{proof}
We first identify all endpoints.  A missing successor can occur
only in the first range of the schedule.  Write its full tape as
$A=Q_x$ and put $y=\operatorname{out}_S(A)=S(x)$.  The current
$B$ tape is its longest valid forward prefix, so
Lemma~\ref{macro:prefix-uniqueness} gives
\begin{equation}\label{macro:sink-identity}
 P(S(y))\ne y,\qquad P(y)=x.
\end{equation}
If $y=0^{n_0}$, then $x=P(0^{n_0})=0^{n_0}$ and
$y=S(0^{n_0})\ne0^{n_0}$, a contradiction.  Thus this designated
vector is a nonzero End-of-Line answer.  Such a vertex has a
predecessor: shorten its forward prefix when its length is positive;
at stage zero use the reverse wraparound copy from the preceding
frame.

A missing predecessor can occur only at a stage $s=c+1$ with
$h\le c<2h$.  Its full tape is $B=Q_y$, where
$y=\operatorname{in}_S(B)$, and its reconstruction prefix has its
maximal length $k_{\mathrm r}(y)<h$.  Hence
\begin{equation}\label{macro:source-identity}
 s=2h-k_{\mathrm r}(y),\qquad S(P(y))\ne y.
\end{equation}
For $y\ne0^{n_0}$ this is again a permitted answer.  The vertex
has a successor obtained by shortening the reconstruction prefix;
if the prefix is empty, the next copying step supplies a successor.
Every proposal in the two copying ranges succeeds in the relevant
direction, except that an incoming reconstruction proposal at stage
$2h$ is covered by the case just considered.  Together with
Lemma~\ref{macro:local-edges}, this exhausts the possible missing
neighbors and proves that no vertex is isolated.

It remains to specify the zero-labelled exception.  The last promise
in \eqref{macro:promises} supplies the unique full transcript
$Q_{0^{n_0}}$, whereas the first two promises give
\[
 S(P(0^{n_0}))=S(0^{n_0})\ne0^{n_0}.
\]
Let $k_*=k_{\mathrm r}(0^{n_0})<h$.  Then $v_*$ is exactly the
vertex at stage $2h-k_*$ whose $B$ tape is $Q_{0^{n_0}}$ and whose
$A$ tape is
$\operatorname{pref}_{\pi_{\mathrm r},k_*}(V_{0^{n_0}})$,
with their prescribed occupations.  Prefix uniqueness proves that
this is the only missing-predecessor vertex with designated vector
zero.  It has no predecessor and has the successor just described.
Evaluate $U_{0^{n_0}}$ and $V_{0^{n_0}}$, locate the first failed
reconstruction terminal comparison, and take the preceding prefix.
This computes $v_*$ and its successor in polynomial time.  This
proves the endpoint claims, including the existence of the public
source.  Other directed paths and directed cycles may remain in the
graph; nothing here requires one connected component.

We next make the local edit interface explicit.  At start frame $c$,
let $j(c)$ be the position in the transition table.  Allot this
position to the content and occupation arrays of the active tape,
and allot the fixed default position $0$ to the two inactive arrays.
Denote these four addresses by $j_D(c)$, for
$D\in\{A,E^A,B,E^B\}$.  If $\mathbf e_{D,j}$ is the unit vector
at position $j$ of array $D$, define, over a field $K$,
\begin{equation}\label{macro:finite-edits}
 h_\sigma(c)=\sum_{D\in\{A,E^A,B,E^B\}}
       \sigma_D\mathbf e_{D,j_D(c)}\in K^{4h},
 \qquad \sigma\in\{-1,0,1\}^4,
\end{equation}
where the integers $-1,0,1$ are mapped into $K$.  For a proposed
Boolean write, choose its content coefficient as the desired bit
minus the old bit and its occupation coefficient as the desired
occupation minus the old occupation.  Set both inactive coefficients
to zero.  Each coefficient lies in $\{-1,0,1\}$, and the resulting
additive edit is precisely the proposal, in either direction.
The vector in \eqref{macro:finite-edits} is public once $c$ and
$\sigma$ are fixed; addition by it is defined on all of $K^{4h}$.
The assertion about valid Boolean proposals does not require a
validity interpretation for other field-valued inputs.

The necessary public data at a frame are its phase, active tape,
write position, rule type, at most two operand addresses, the old
content and occupation addresses, and, when applicable, the extra
comparison address.  A forward fill uses its forward rule; a reverse
fill uses its reconstruction rule.  Lemma~\ref{macro:local-edges}
shows that the existence flag is the new comparison result in the
only direction where failure is possible, and is one otherwise.
These are fixed Boolean decisions once the finitely many addressed
bits and public rule data are supplied.  Define the decisions on
arbitrary Boolean reads by the same formulas, using a fixed default
rule and address for unused roles.  Their agreement with graph
neighbors is asserted on valid input vertices.  No global validity
check or full circuit evaluation is hidden in an existence flag.

For completeness, the bounded-multiplicity assertion includes more
than internal gate edges.  Each internal value has boundedly many
uses after copy-tree normalization.  Each designated output is used
once by the other circuit's input-copy rule and participates in only
a bounded number of terminal or external-vector comparisons.  Each
external-vector bit is used once as an input bit and once for its
corresponding terminal comparison.  Copied-prefix equations use each
full-tape position once.  Occupation indicators condition only the
rules at their positions; prefix monotonicity, if expressed by
adjacent-position tests, reads a position at most twice in each
order.  Thus all local content, bitness, occupation, gate, copy,
comparison, and dummy conditions use bounded-arity records, and each
logical position appears as an operand only boundedly often across
the fixed list of phases.  The requirement that the occupied prefix
have its specified length is separate: one may either inspect the
prescribed occupation pattern directly, or use prefix monotonicity
together with the occupation sum.  The latter is a single global
count condition for each tape and will be represented separately in
the algebraic construction.

All normalization and ordering operations are deterministic.  Their
records and the two permutations occupy $O(h\log(N+2))$ bits.
Tabulate the transition data for the $4h$ frames in the same space
bound, including both directions and defaults.  Scanning the arrays
and these records tests membership in $\mathcal V_s$ in polynomial
time.  The designated endpoint vectors and the circuit inequalities
\eqref{macro:sink-identity} and \eqref{macro:source-identity} give a
polynomial-time endpoint decoder.  At every fixed stage a valid
vertex is determined by the input of its full transcript, since
its other tape is a unique computed or copied prefix.  Consequently
the graph has at most $4h\,2^{n_0}$ vertices, even though an individual
label has only $O(N)$ bits.  The construction outputs the normalized
records, public tables, and algorithms, never this potentially
exponential vertex list.  All constants in these bounds depend only
on the fixed Boolean basis and padding cutoff.  This proves the
remaining uniformity and size assertions.
\end{proof}
